% Overleaf: compiler pdfLaTeX. This file needs no images or other project files.
\documentclass[12pt,a4paper]{article}
\usepackage[T2A]{fontenc}
\usepackage[utf8]{inputenc}
\usepackage[russian]{babel}
\usepackage[margin=2.3cm]{geometry}
\usepackage{enumitem}
\usepackage[unicode,hidelinks]{hyperref}
\setlength{\parindent}{1.25cm}
\setlength{\parskip}{0pt}
\setlist[enumerate]{leftmargin=*,topsep=4pt,itemsep=3pt,parsep=0pt}
\emergencystretch=2em

\begin{document}

\section{Что мы сделали}

Мы собрали на Concordia~2.4.0 прототип изменения коллективных предпочтений.
Шесть вымышленных участников выбирают, что развивать в библиотеке:
цифровой доступ или печатные книги. У каждого есть свои потребности и бюджет.
Агенты получают новые сведения, читают аргументы собеседников и объясняют
своё решение.

В новой версии собеседники меняются между раундами без повторных встреч.
Мы подключили bonsai2-27b через API и GPT-6-Luna через Codex.
Промпты, ответы и условия опытов сохранены. Есть HTML-демонстрация,
которая показывает записанные результаты без новых обращений к моделям.

\section{Как устроена Concordia}

Агент собирается из компонентов. Мы использовали \texttt{EntityAgent},
компоненты памяти и наблюдений, а также \texttt{ConcatActComponent}
для сборки контекста и вызова модели. Это позволяет отдельно менять
профиль, память и способ принятия решения.

Наблюдения записываются в \texttt{ListMemory}. Перед действием мы выбираем,
что попадёт в промпт. В текущей версии это официальные сведения,
последняя реплика выбранного партнёра и собственная история аргументов.
Хранение записи и её использование в следующем ответе --- разные операции.

Среду и расписание мы написали сами; стандартный Game Master не запускали.
Среда задаёт новости и пары участников. Агент видит сообщение партнёра
из предыдущего раунда, затем все новые ответы фиксируются вместе.
Это один обмен позициями за раунд, а не полноценный диалог внутри пары.

\section{Плюсы и минусы для нашей задачи}

\textbf{Плюсы.} Компоненты удобно заменять по отдельности.
Через общий интерфейс модели удалось подключить и API, и Codex,
а для проверки инфраструктуры использовать численную заглушку.
Состояние агента можно сохранять; восстановление мы проверили в первом
протоколе. Такой подход подходит для небольших управляемых экспериментов.

\textbf{Минусы.} Для нашего контролируемого сценария пришлось отдельно
сделать расписание, доставку сообщений, журналы и оценку результатов.
Содержимое памяти нужно подбирать явно: лишняя история увеличивает
контекст и может закреплять старые ошибки. Качество рассуждений всё равно
зависит от модели и промпта. Concordia сама по себе эту проблему не решает.
В нашем прототипе пока нет динамического доверия и отдельной памяти отношений.

\section{Что проверили и что получилось}

Мы сравнили новости и общение по отдельности, разные исходные профили,
варианты памяти, постоянных и меняющихся собеседников. Проверили адресную
доставку сообщений и независимость результата от порядка вызова агентов.

В первом сценарии агенты часто застывали после первого раунда
и повторяли общие доводы. При других исходных профилях эффект общения
менял направление. Поэтому точное повторение одного запуска ещё
не означает устойчивости результата.

Новый сценарий сделал ошибки заметнее. У bonsai встречались объяснения,
где стоимость~80 называлась меньшей, чем~40, а стоимость~120 ---
укладывающейся в бюджет~100. При этом факты в отдельных полях ответа
были записаны правильно. Значит, проверки формата и извлечения фактов
недостаточно: нужно проверять сам довод.

В просмотренных ответах Luna последовательнее учитывала бюджет
и исправление ошибочной стоимости. Но это предварительное сравнение:
у Codex свои системные инструкции и настройки, а независимой разметки
всех ответов пока нет.

Смена пар работает технически, однако у bonsai численные траектории
с новыми и постоянными собеседниками совпали. Без собственной истории
в одном условии больше агентов пересмотрели мнение. Добавление памяти
не всегда улучшает результат. Социальная правдоподобность пока не проверена.

\section{Что делаем дальше}

\begin{enumerate}
    \item Задать критерии качества аргумента: ошибки в бюджете,
    противоречия фактам, выдуманные обстоятельства и повторение доводов.
    Разметить ответы обеих моделей по одним правилам.
    \item Повторить опыты с другими профилями и несколькими генерациями.
    Для сравнения моделей по возможности выровнять способ подключения
    и параметры запросов.
    \item Отдельно проверить влияние собеседника: сравнить общение,
    отсутствие сообщений и перемешанные реплики при одинаковых новостях.
    \item Сравнить полную историю с кратким резюме и памятью,
    которая явно учитывает исправления. После этого выделить в свою
    библиотеку интерфейсы агента, памяти, событий, контактов и модели.
\end{enumerate}

Минимальный прототип работает. Следующая задача --- проверить,
что агенты действительно учитывают новые сведения и чужие аргументы,
а не только воспроизводят удачно сформулированный промпт.

\smallskip
\noindent Код и журналы опытов:
\url{https://github.com/QwichaFicha/llm-opinion-dynamics}.

\end{document}
